\documentclass[12pt,a4paper]{article}

% --- Encodage et langue
\usepackage[utf8]{inputenc}
\usepackage[T1]{fontenc}
\usepackage[french]{babel}

% --- Police Times New Roman
\usepackage{mathptmx}

% --- Mise en page
\usepackage[top=2.5cm, bottom=2.5cm, left=3cm, right=2.5cm]{geometry}
\usepackage{setspace}
\onehalfspacing

% --- Justification et typographie
\usepackage{microtype}
\usepackage{parskip}
\setlength{\parindent}{1.25cm}
\setlength{\parskip}{0pt}

% --- Liens sans couleur
\usepackage[colorlinks=false, pdfborder={0 0 0}]{hyperref}

% --- Tableaux
\usepackage{booktabs}
\usepackage{array}
\usepackage{tabularx}

% --- Pas de numérotation des sections
\setcounter{secnumdepth}{0}

% --- En-tête et pied de page
\usepackage{fancyhdr}
\pagestyle{fancy}
\fancyhf{}
\fancyhead[L]{\small Groupe 1 -- ERI-ESI 2017}
\fancyhead[R]{\small Traitement de données}
\fancyfoot[C]{\thepage}
\renewcommand{\headrulewidth}{0.4pt}

% ============================================================
\begin{document}

% --- Page de garde
\begin{titlepage}
\centering
\vspace*{1.5cm}

{\large \textbf{École Nationale de la Statistique\\et de l'Analyse Économique}}\\[0.3cm]
{\normalsize Dakar, Sénégal}

\vspace{2cm}
\rule{\linewidth}{0.5pt}\\[0.4cm]
{\LARGE \textbf{Traitement de données\\[0.3cm]
Enquête Régionale Intégrée sur l'Emploi\\
et le Secteur Informel (ERI-ESI 2017)}}\\[0.4cm]
\rule{\linewidth}{0.5pt}

\vspace{1.2cm}
{\large \textit{Rapport de synthèse -- Groupe 1}}

\vspace{2cm}
\begin{tabular}{ll}
\textbf{Membres du groupe :} & Ndeye Aïssatou CISSE \\
                              & Papa Maguette DIOP \\
                              & Aïssatou GUEYE \\
                              & Awa GUEYE \\
                              & Marc MARE \\[0.3cm]
\textbf{Qualité :}            & Analystes statisticiens en fin de formation \\[0.3cm]
\textbf{Dépôt GitHub :}       & \href{https://github.com/papamagattediop/Traitements\_de\_donnees-ERI-ESI-}%
                                {\texttt{papamagattediop/Traitements\_de\_donnees-ERI-ESI-}} \\
\end{tabular}

\vfill
{\normalsize Année académique 2024-2025}
\end{titlepage}

% ============================================================
\section*{1.\quad Contexte et objectifs}

L'ERI-ESI (Enquête Régionale Intégrée sur l'Emploi et le Secteur Informel) est une
opération statistique conjointe conduite en 2017 par l'ANSD, l'AFRISTAT et l'UEMOA
dans plusieurs pays membres. Le volet sénégalais, objet de ce travail, porte sur
l'emploi, les conditions d'activité et le secteur informel.

L'objectif du projet est de construire un pipeline de traitement reproductible à
partir des fichiers bruts de l'enquête, en produisant deux tables analytiques
consolidées (individus et ménages) ainsi qu'un ensemble de contrôles de qualité et
d'estimations de référence. Le pipeline est conçu pour être généralisable aux autres
volets pays de l'ERI-ESI, moyennant l'adaptation d'un bloc de configuration centralisé.

% ============================================================
\section*{2.\quad Sources de données}

Deux fichiers \texttt{.dta} constituent les entrées du pipeline.

\begin{table}[h!]
\centering
\begin{tabularx}{\linewidth}{lXrr}
\toprule
\textbf{Fichier} & \textbf{Contenu} & \textbf{Obs.} & \textbf{Variables} \\
\midrule
\texttt{individus\_TP.dta} & Roster des membres du ménage & 120 693 & 53 \\
\texttt{emploi\_Tp.dta}    & Module emploi individuel      & 120 689 & 492 \\
\bottomrule
\end{tabularx}
\caption{Fichiers sources (volet Sénégal, ERI-ESI 2017)}
\end{table}

Les deux fichiers sont appariés sur la clé \texttt{hh1 + hh2 + m1}
(grappe, ménage, numéro d'ordre de l'individu). Quatre individus du roster
n'ont pas de correspondance dans le module emploi, ce qui est cohérent avec
le champ d'application de ce module, restreint aux personnes de dix ans et plus.
La base fusionnée regroupe ainsi 120 693 lignes et 543 variables.

Les fichiers de données, trop volumineux pour être versionnés (807 Mo, 82 Mo
et 99 Mo respectivement), sont mis à disposition par voie externe, conformément
aux instructions précisées dans le fichier \texttt{README.md} du dépôt.

% ============================================================
\section*{3.\quad Pipeline de traitement}

Le pipeline est entièrement implémenté en Stata 18. Il s'articule en six scripts
exécutés séquentiellement depuis la racine du dépôt.

\begin{table}[h!]
\centering
\begin{tabularx}{\linewidth}{llX}
\toprule
\textbf{Script} & \textbf{Entrée} & \textbf{Sortie} \\
\midrule
\texttt{01\_fusion.do}
  & \texttt{individus\_TP.dta}, \texttt{emploi\_Tp.dta}
  & \texttt{base\_individus\_emploi\_fusionnee.dta} \\[4pt]
\texttt{03\_demographie\_geo.do}
  & Base fusionnée
  & \texttt{individus\_demo\_geo.dta} \\[4pt]
\texttt{03\_emploi\_principal.do}
  & Base fusionnée
  & \texttt{emploi\_principal.dta} \\[4pt]
\texttt{04\_emploi\_secondaire.do}
  & Base fusionnée
  & \texttt{emploi\_secondaire.dta} \\[4pt]
\texttt{05\_menages.do}
  & \texttt{individus\_demo\_geo.dta}
  & \texttt{menages.dta} \\[4pt]
\texttt{06\_consolidation.do}
  & Sorties des modules 03 à 05
  & \texttt{individus\_consolide.dta}, \texttt{menages\_consolide.dta} \\
\bottomrule
\end{tabularx}
\caption{Séquence d'exécution du pipeline}
\end{table}

La configuration pays (chemins, noms de fichiers, clés de fusion) est centralisée
dans \texttt{config.do}. L'adaptation à un autre pays de l'UEMOA ne requiert que
la modification de ce bloc, sans toucher à la logique de traitement.

% ============================================================
\section*{4.\quad Variables construites}

\subsection*{4.1\quad Module démographie et géographie}

Ce module construit les variables individuelles de base à partir du roster :
sexe, âge et groupe d'âge (huit tranches), lien de parenté avec le chef de
ménage (détaillé et regroupé en cinq postes), situation matrimoniale, niveau
d'études consolidé (six niveaux) et branche d'études. La région et le milieu
de résidence (urbain/rural) sont également retenus. Cinq indicateurs de
contrôle (flags) signalent les incohérences détectées sans modifier les
observations.

\subsection*{4.2\quad Module emploi principal}

La situation d'activité est issue de la variable synthétique \texttt{sitac}
produite par l'ANSD, regroupée en quatre postes : actif occupé, chômeur BIT,
main-d'oeuvre potentielle et inactif. Pour les actifs occupés, le module
construit le statut dans l'emploi, le type (salarié/indépendant), le secteur
d'activité selon la nomenclature ISIC Rév. 4 (21 sections et regroupement en
quatre familles), les heures travaillées par semaine, le sous-emploi lié à la
durée, le revenu mensuel consolidé et les indicateurs de formalité de l'unité
de production et de l'emploi.

Cinq correspondances erronées ont été identifiées lors de la vérification du
dictionnaire initial par croisement avec le questionnaire et les métadonnées
du fichier. Les corrections apportées figurent dans le tableau ci-dessous.

\begin{table}[h!]
\centering
\begin{tabularx}{\linewidth}{lllX}
\toprule
\textbf{Variable} & \textbf{Initiale} & \textbf{Corrigée} & \textbf{Note} \\
\midrule
Situation matrimoniale    & M5     & M25    & M5 est le lieu de naissance \\
Secteur ISIC principal    & AP4    & AP2\_Niv1 & AP4 est le secteur institutionnel \\
Heures travaillées        & AP6A   & AP10C  & AP6A est le régime fiscal \\
Revenu principal          & AP8A1  & AP13A/B/C & AP8A1 est l'ancienneté dans l'emploi \\
Raison d'inactivité       & SE9    & SE11   & SE9 est la disponibilité à travailler \\
\bottomrule
\end{tabularx}
\caption{Corrections apportées au dictionnaire de variables}
\end{table}

\subsection*{4.3\quad Module emploi secondaire}

Le module traite le premier emploi secondaire déclaré (suffixe \texttt{\_AS1}),
seul statistiquement significatif (118 répondants seulement pour le second).
Il construit le statut, le type d'emploi, le secteur d'activité, les heures
hebdomadaires (converties depuis des minutes par jour renseignées dans
\texttt{AS9C}), le revenu mensuel et la formalité de l'emploi. Une variable
indicatrice signale les codes sectoriels potentiellement issus d'une révision
antérieure de la nomenclature ISIC.

\subsection*{4.4\quad Table ménages}

La table ménages est construite à partir du chef de ménage identifié dans le
roster. Elle contient les caractéristiques démographiques et géographiques du
chef, la structure du ménage (taille, répartition par sexe, nombre de membres
de moins de 15 ans et de 15 ans et plus) ainsi que, après consolidation, les
agrégats d'emploi au niveau ménage (nombre d'actifs occupés, de salariés,
d'emplois informels et ratio de dépendance).

% ============================================================
\section*{5.\quad Contrôle qualité}

Chaque module produit deux fichiers de sortie complémentaires : un fichier de
contrôles de cohérence (\texttt{qc\_*.csv}) listant les anomalies détectées avec
leur effectif, et un fichier d'estimations pondérées (\texttt{estimations\_*.csv})
permettant la comparaison avec les chiffres publiés par l'ANSD. Le module de
consolidation produit également des contrôles croisés portant sur la cohérence
entre les tables individus et ménages.

Les estimations clés obtenues sur le volet emploi, comparées aux valeurs de
référence du rapport final ANSD (2017), sont présentées ci-dessous à titre
indicatif.

\begin{table}[h!]
\centering
\begin{tabularx}{\linewidth}{Xrr}
\toprule
\textbf{Indicateur} & \textbf{Pipeline} & \textbf{ANSD} \\
\midrule
Taux de chômage BIT (\%)                    & recalculé & 2,9 \\
Taux de salarisation (\%)                   & recalculé & 38,6 \\
Part des indépendants parmi les occupés (\%) & recalculé & 66,1 \\
Taux d'emploi informel (\%)                 & recalculé & 95,4 \\
Secteur primaire parmi les occupés (\%)     & recalculé & 24,7 \\
Secteur commerce parmi les occupés (\%)     & recalculé & 27,6 \\
\bottomrule
\end{tabularx}
\caption{Indicateurs d'emploi -- valeurs de référence ANSD (2017)}
\end{table}

La colonne \og Pipeline \fg{} est à renseigner après exécution complète sur les
données réelles, les fichiers bruts n'étant pas versionnés dans le dépôt.

% ============================================================
\section*{6.\quad Reproductibilité et généralisation}

L'ensemble du pipeline est reproductible depuis la racine du dépôt en exécutant
les six scripts dans l'ordre indiqué au tableau 2. La seule condition préalable
est la présence des fichiers sources dans le répertoire \texttt{input/SEN/}.
L'adaptation à un autre pays de l'UEMOA ne nécessite que la modification du
bloc global dans \texttt{config.do} (code pays, noms de fichiers, clés de
fusion et noms de variables si la nomenclature du questionnaire diffère).

Le dépôt est organisé comme suit : \texttt{scripts/stata/} pour les scripts de
traitement, \texttt{docs/} pour la documentation et le présent rapport,
\texttt{input/} et \texttt{output/} pour les données (non versionnées). Un
fichier \texttt{README.md} décrit l'organisation du projet et les étapes
d'exécution.

% ============================================================
\section*{7.\quad Conclusion}

Ce travail produit un pipeline de traitement structuré et documenté pour les
données ERI-ESI 2017, volet Sénégal. Les deux tables analytiques finales
(individus et ménages) couvrent l'ensemble des dimensions thématiques
accessibles dans les fichiers du volet emploi : démographie, géographie,
situation d'activité, emploi principal et secondaire, formalité et revenus.
La démarche de contrôle qualité intégrée à chaque module garantit la traçabilité
des choix méthodologiques et facilite la comparaison avec les statistiques
officielles publiées par l'ANSD.

\vspace{1.5cm}
\noindent Nous tenions à vous remercier pour l'enseignement dispensé tout au long
de ce cours et restons à votre disposition pour toute précision.

\vspace{1cm}
\noindent\textit{Cordialement,}

\vspace{0.5cm}
\noindent\textbf{Membres du Groupe 1}\\
Ndeye Aïssatou CISSE, Papa Maguette DIOP, Aïssatou GUEYE, Awa GUEYE, Marc MARE\\
\textit{Analystes statisticiens en fin de formation}

\end{document}
